\section{与现有工具的对比分析}
\label{sec:comparison}

\subsection{对比范围与方法}

本章将 e2m2e 置于现有航天软件生态中定位。对比对象共 10 个，覆盖四类
形态：通用任务分析平台（GMAT、STK、ATK）、开源天体力学库（Orekit、
pykep、poliastro）、自建路线的起点与精度标尺（SciPy）、生态互补件
（r2s2、Cesium、ootk）。调研于 2026 年 8 月完成，版本锚点为
e2m2e 5.8.10：对 GMAT、Orekit、pykep、poliastro、ootk、r2s2 做了
本地代码级调研（逐仓库 grep 与源码通读），对商业闭源的 STK、ATK 及
纯渲染层的 Cesium 采用官方文档与第三方文献调研；文中所有数字均可在
各对象的仓库或官方文档中复核。

为保证叙述公允，本章遵循三条原则：

\begin{enumerate}
  \item \textbf{无对拍不倍数}：凡缺乏同场景对拍数据的对象
        （STK、Orekit、ATK），只作架构性对比，不宣称性能倍数；
  \item \textbf{性能分区陈述}：对 SciPy 自建路线的加速可定量
        （回调开销是结构性的）；对 GMAT 的差异在并行架构层面陈述；
  \item \textbf{精度谈对齐不谈胜负}：e2m2e 的力模型层以 GMAT 数学规格
        为移植基准，精度关系是``对齐声明 + 组件级移植一致性''，
        而非``谁更准''。
\end{enumerate}

\subsection{生态全景}

\begin{table}[htbp]
\centering\footnotesize
\caption{对比对象全景（2026 年 8 月版本锚点）}
\label{tab:landscape}
\begin{tabularx}{\textwidth}{@{}llXX@{}}
\toprule
\textbf{工具} & \textbf{开发方 / 许可} & \textbf{定位（语言）} & \textbf{与 e2m2e 的关系} \\
\midrule
GMAT & NASA / Apache-2.0 & 通用任务分析平台，GUI + 插件（C++，约 85 万行） &
e2m2e 的力模型对齐基准；无 CR3BP 内核 \\
STK & Ansys / 商业闭源 & 数字任务工程全功能平台 &
行业事实标准；Astrogator 有 CR3BP 框架但无族自动化 \\
ATK & 国防科大 / 闭源二进制 & STK 国产替代，300+ 单位 &
通用 MCS 框架；公开材料未见三体专用工具 \\
Orekit & CS GROUP / Apache-2.0 & 飞行动力学基础设施库（Java，约 47 万行） &
最成熟基础设施；三体仅 Halo/Lyapunov 两型构件 \\
pykep & ESA / MPL-2.0 & 行星际轨迹优化研究库（C++23 + pybind） &
三体仅积分器与基准题；非运行级工具 \\
poliastro & 社区 / MIT & 教学级天体力学（Python + numba） &
\textbf{已归档}（2023-10）；三体停在 contrib 草稿 \\
SciPy & 社区 / BSD 系 & 通用科学计算（Python + C/Fortran） &
自建路线的起点，同时是 e2m2e 的精度标尺 \\
ootk & thkruz / AGPL-3.0 & Web 端地球轨道 SSA 工具库（TypeScript） &
赛道不同，几乎零重叠，互补 \\
r2s2 & 紫金山天文台 / GPL-3.0 & 相对论时空坐标转换（Python，约 1.2 千行） &
\textbf{e2m2e 的底层依赖}，非竞品 \\
Cesium & Cesium GS / Apache-2.0 & 三维可视化引擎（JS/TS，WebGL） &
纯渲染层零计算，天然前后端组合 \\
\bottomrule
\end{tabularx}
\end{table}

\subsection{核心差异化轴：地月三体设计自动化}
\label{subsec:three-body-axis}

e2m2e 与所有对比对象的分水岭在于：三体能力是``设计流水线''还是
``零件或框架''（表~\ref{tab:threebody}）。

\begin{table}[htbp]
\centering\small
\caption{地月三体能力对比（白皮书核心差异化轴）}
\label{tab:threebody}
\begin{tabularx}{\textwidth}{@{}lXX@{}}
\toprule
\textbf{工具} & \textbf{CR3BP 动力学 / 族生成} & \textbf{低能转移 / 保持 / 星历修正} \\
\midrule
\textbf{e2m2e} & \textbf{原生}三模型（CR3BP/BCR4BP/星历 $N$ 体）互转；
\textbf{八族自动生成} + 延拓 + 能量窗口批量 &
流形拼接 + LGA/WSB；三控制律 + Monte Carlo；
初猜 $\to$ 星历修正（20 m 收敛）$\to$ 高保真三段链 \\
GMAT & 无 CR3BP 模块（全仓 grep 零命中）；平动点仅为几何点，旋转系靠
用户自拼 & 均无；星历传播为回放器形态 \\
STK & Astrogator 11.7 起内建 CR3BP 力模型 + 差分打靶 &
无模板化族生成（未见公开证据）；低能转移靠 MCS 自建；无保持控制律 \\
ATK & 高精度 $N$ 体数值传播 + 通用 MCS（11 段型 $\times$ 60 约束）&
公开材料未宣称族生成/流形；有复现 GMAT 算例的 Halo 维持论文证据 \\
Orekit & CR3BPForceModel + STM，构件式 &
仅 Halo/Lyapunov 两型；流形仅初值单点；无保持；打靶可配星历（用户装配） \\
pykep & heyoka 符号工厂 + 42 维变分方程 &
无族生成（grep 零命中）；无流形；14 个地月低推力\textbf{基准题}——
其输入恰是 e2m2e 的产出物 \\
其余四者 & SciPy 零轨道知识；poliastro 三体停在 contrib 草稿；
ootk/Cesium 无三体代码 & --- \\
\bottomrule
\end{tabularx}
\end{table}

结论是清晰的：\textbf{八个对比对象无一同时具备``八族生成 + 流形低能
转移 + 三体轨道保持 + 星历修正链''}。GMAT 的三体是几何点加 $N$ 体
求和，STK 提供框架但工作流靠手工搭序列，Orekit 提供构件但装配靠用户，
pykep 把三体当作优化载体而非设计对象，poliastro 的三体从未进入主库。
``地月三体设计自动化''因此是 e2m2e 的独特生态位。

\subsection{代表性对象逐评}

\subsubsection{通用任务分析平台：GMAT、STK、ATK}

\textbf{GMAT}（NASA，约 85 万行 C++）是与 e2m2e 关系最``近''的对象：
e2m2e 的 PD78/RK89 权重、锥形阴影、相对论三项、STM 步进误差控制语义、
EOP/闰秒 fixtures 均以 GMAT R2026a 数学规格为移植基准，力模型层可以
说是``GMAT 数学规格的 Rust 重实现 + 扩展''。反过来，GMAT 的求解器与
优化生态（DifferentialCorrector、Yukon SQP、CSALT 直接配点、Bulirsch--
Stoer 积分器）与文档级 V\&V 体系（数学规格书、飞行资格文档）是 e2m2e
尚未有的厚度；其无并行传播内核（物理循环 grep 零命中）则是 2000 年代
桌面软件形态与 e2m2e 服务化形态的代际差异。二者是``对标--差异化''
关系：GMAT 出校核与运营分析，e2m2e 出三体设计流水线。

\textbf{STK}（Ansys，商业闭源）是行业事实标准，Astrogator 自 11.7 起
内建 CR3BP 力模型与差分打靶，Artemis II 任务的 Orion 跟踪可视化即由
STK 驱动。其与 e2m2e 的对比可归为四轴：许可模式（闭源商业 vs
Apache-2.0）、三体自动化（框架 vs 一键族生成）、成本（政府采购记录中
单 Pro token 订单 11.45 万美元 vs 免费）、可编程性（Integration 许可
模块 vs 原生 MCP）。STK 的覆盖/链路/调度/定轨（ODTK）模块生态 e2m2e
均不覆盖，两者在典型工作流中是``开源计算 + 商业呈现''的分层组合。

\textbf{ATK}（国防科大，闭源二进制）走``中国版 STK''路线：积木式
机动规划框架（11 种轨道段类型、19 种停止条件、11 类 60 种约束组件，
统一 NLP + 微分修正/SQP/智能优化三类求解器），有天宫、神舟、天舟等
载人航天实战背书。其论文算例复现了地月自由返回与 GMAT 内置的 L2 Halo
维持案例（四次机动 $-0.25/1.64/8.05/6.44$ m/s），说明框架可承载三体
问题；但公开材料中未见 CR3BP 传播器、平动点轨道族生成或不变流形
自动化——这仍是 e2m2e 最干净的差异点。``平台 vs 库''、``闭源 vs
开源''、``通用 MCS vs 三体专用流水线''是两者对比的三条主线。

\subsubsection{开源天体力学库：Orekit、pykep、poliastro}

\textbf{Orekit}（CS GROUP，约 47 万行 Java，2002 年起 24 年积累）是
工程级基础设施的标杆：帧/时间链对官方 SOFA 库逐元素对照至 $10^{-12}$，
传播器谱系覆盖数值/解析/DSST 半解析/SGP4/GNSS，定轨全家桶（批最小
二乘 + EKF/UKF）、80 余种事件检测器都是 e2m2e 没有的深度。但其三体
能力是构件式的——动力学模型 + Richardson 初猜 + 打靶零件齐备，
轨道类型仅 Halo/Lyapunov 两枚举，无族延拓、无 Lissajous/DRO、无流形
批量积分与拼接、无保持控制律。两者同为 Apache-2.0，``Orekit 管地球段
与定轨、e2m2e 管地月段设计''是自然的分工。

\textbf{pykep}（ESA，kep3，MPL-2.0）是行星际轨迹优化研究库：heyoka
LLVM JIT Taylor 积分器（容差至 $10^{-16}$）+ pagmo 进化优化，MGA
编码、Sims--Flanagan/ZOH 低推力转录、Pontryagin 间接法，官方自述
``非运行级飞行动力学工具''。其三体能力止于符号动力学工厂与 14 个
地月低推力基准题（TOPS：Halo-to-Halo、DRO-to-NRHO 等）——未实现
平动点计算、族构造、单值矩阵与流形打靶；而 TOPS 基准题的输入
（Halo/DRO/NRHO 标称轨道）恰是 e2m2e 的输出。``e2m2e 生成标称轨道
$\to$ pykep + pygmo 做低推力全局优化 $\to$ e2m2e 高保真验证与
保持''是一条现成的分工链。

\textbf{poliastro}（MIT，2023 年 10 月归档）对本白皮书的意义在于
\textbf{生态空位的证据}：它的单位安全 API 与 plotly 可视化是教学级
精品，但地月三体能力停留在 contrib 草稿目录（其 README 自列待办：
微分修正、族延拓、Halo 三阶解析解、单值矩阵、流形均未实现），力模型
仅 6 个摄动函数、测试容差 $10^{-4}\sim10^{-2}$ 级。Python 生态中最
活跃的天体力学教学库归档且未竟三体之作，正是 e2m2e 所填的空。

\subsubsection{自建路线与精度标尺：SciPy}

SciPy 对轨道一无所知——无时间系统、无参考架、无星历、无力模型，
但 \texttt{solve\_ivp} + \texttt{root} 搭 CR3BP 打靶是学术界自建的
事实标准路线，也是 e2m2e 用户会先尝试的``下位替代''。性能上其显式
RK 族为 Python 实现，每步回调右端函数，紧容差大批量打靶场景比编译
内核慢 1--3 个数量级（社区共识；e2m2e 侧对应证据：同一线程化问题
Python 多进程仅得 Rust 下沉后加速的约 $1/10$）。精度上两者打平——
e2m2e 的测试体系恰恰以 scipy DOP853（$10^{-13}$）为金标准参考解，
LEO+J2 一天弧一致性 $<10^{-9}$。差距不在精度，而在``要达到同样精度，
用户需自建整个领域层''。

\subsubsection{生态互补件：r2s2、Cesium、ootk}

\textbf{r2s2}（紫金山天文台，约 1150 行纯 Python）不是竞品而是
e2m2e 的底层依赖：GCRS 与 EBCRS 之间的相对论时空转换即由它
承载，反向转换默认收敛 1 ns / 1 mm。e2m2e 对它的独立验证（往返一致
性、与 SPICE 链路差分、与 ERFA \texttt{dtdb} 差分）反而构成其最主要的
测试覆盖。许可上 r2s2 为 GPL-3.0，闭源集成方需注意传染性。

\textbf{Cesium}（CesiumJS，Apache-2.0）是纯渲染层，零轨道力学——
卫星位置必须外部喂入，历史上本就``为跟踪空间物体而生''，对地月大
尺度场景有一等公民支持（ICRF 惯性系渲染、Artemis II 环月跟踪实证）。
e2m2e 的可视化恰是短板，``EphemerisTable $\to$ CZML $\to$ Cesium''
是零许可冲突的组合路径（Orekit + OreCZML 已验证同一模板）；需牢记
的分工是：Cesium 忠实渲染喂给它的轨迹，轨迹本身的正确性由 e2m2e
负责。

\textbf{ootk}（TypeScript，AGPL-3.0）在另一条赛道：SGP4/TLE 编目、
传感器、碰撞、Web 端 SSA，与 e2m2e 几乎零重叠。其价值有二：一是
互补（编目/观测端归 ootk，高保真任务设计归 e2m2e）；二是方法论
参照——ootk 以 AFSPC 官方 SGP4Prop 位级快照与 Vallado 验证向量做
回归、以 SLR 精密星历做独立真值，这种``对外部权威实现的位级对齐''
正是 e2m2e 验证体系升级可借鉴的方向。

\subsection{性能与精度：证据的分区陈述}

\begin{table}[htbp]
\centering\small
\caption{并行/加速架构与公开量化数字}
\label{tab:perf}
\begin{tabularx}{\textwidth}{@{}lXl@{}}
\toprule
\textbf{工具} & \textbf{架构} & \textbf{公开数字} \\
\midrule
\textbf{e2m2e} & Rust 内核 + Rayon（打靶/网格/流形/WSB） &
可复现，benchmark 脚本在库 \\
GMAT & 无并行传播内核（单实例顺序仿真） & 无 \\
STK & Parallel Computing 多核分发 & 8 核 180 样本 70 s（仅此） \\
Orekit & PropagatorsParallelizer & 无公开数字 \\
pykep & heyoka LLVM JIT + 按容差缓存 & SF 求值 5.25 $\mu$s；
``千条轨迹/秒'' \\
SciPy & 单线程，Python 右端回调 & 比编译内核慢 1--3 个量级 \\
ootk & 纯 JS 单线程（轻计算 $\times$ 海量调用） & SGP4 391 ns/次 \\
\bottomrule
\end{tabularx}
\end{table}

e2m2e 自身的性能数字均可复现：转移网格搜索 Rust 化加速
5.5--10.4$\times$；180 天 Halo 端到端设计 223 s $\to$ 63 s；SRP 走
Rust STM 快速路径后打靶提速 9.4$\times$；低推力解析雅可比 3--24$\times$。
按本章开头的原则，这些数字只与 SciPy 自建路线对量级、与 GMAT 谈架构，
与 STK/Orekit/ATK 不作相互倍数声明——三者均无同场景公开对拍数据。

精度方面（表~\ref{tab:prec}），各家的锚点与验证哲学不同：Orekit 与
ootk 的``对权威参考逐元素/位级对照''证据链最完整；pykep 在无量纲
归一化单位下达到 $10^{-13}$；e2m2e 的积分器对解析解 $10^{-9}\sim
10^{-11}$、坐标链 $<10^{-12}$、力模型对齐 DFH 至亚百米（Halo 7 天
88 m）——但按 ADR 0013 的``按定义验证''原则，外部对拍已撤出仓库，
 GMAT 侧仅保留可自跑的对比脚本。精度不是各家分化点，可解问题域才是。

\begin{table}[htbp]
\centering\small
\caption{精度验证锚点}
\label{tab:prec}
\begin{tabularx}{\textwidth}{@{}lXl@{}}
\toprule
\textbf{工具} & \textbf{外部真值锚点} & \textbf{代表性数字} \\
\midrule
\textbf{e2m2e} & DFH / GMAT fixtures / scipy DOP853 / 解析解 &
Halo 7 天 88 m；积分器 $10^{-9}\sim10^{-11}$；ITRF $<10^{-12}$ \\
GMAT & V\&V 文档体系 + 机构背书 & 默认步容差 $10^{-11}$ \\
STK & 第三方论文 / 社区交叉 & RKF78 步容差 $10^{-13}$ 可配 \\
ATK & 论文算例（含 GMAT Halo 算例复现） & 交会 100 m 瞄准达 100.23 m \\
Orekit & 官方 SOFA 库逐元素对照 & 帧变换 $10^{-12}$ \\
pykep & 硬编码高精度真值 & CR3BP $\le 10^{-13}$（容差 $10^{-16}$） \\
ootk & AFSPC 位级快照 + Vallado 向量 + SLR 星历 & SGP4 $5\times10^{-9}$ km \\
r2s2 & （由依赖方 e2m2e 补测） & 迭代收敛 1 ns / 1 mm \\
\bottomrule
\end{tabularx}
\end{table}

\subsection{组合生态：地月设计引擎}

把视角从竞争切换到组合，e2m2e 在生态图中的位置是\textbf{地月设计
引擎}：上游有时空与星历供给（NASA SPICE/DE 历表、r2s2 相对论时空
转换），下游有现成的开源消费件（图~\ref{fig:eco}）——Cesium 渲染、
GMAT 校核与运营分析、Orekit 承接地球段与定轨、pykep 接手低推力
全局优化（其地月基准题的输入正是 e2m2e 的产出）。每一条组合链的
许可均可共存（r2s2 的 GPL-3.0 需闭源集成方注意）。

\begin{figure}[htbp]
\centering
\begin{tikzpicture}[
  font=\footnotesize,
  box/.style={draw=e2gray, rounded corners=1.2mm, fill=white,
    minimum height=0.75cm, align=center, inner sep=4pt},
  core/.style={draw=e2blue, very thick, rounded corners=1.5mm, fill=e2light,
    minimum height=1.5cm, minimum width=3.6cm, align=center, font=\small},
  arr/.style={-{Stealth[length=2.2mm]}, thick, e2blue},
]
\node[core] (e2) at (0,0) {\textbf{e2m2e}\\地月设计引擎\\族生成 $\cdot$ 转移 $\cdot$ 保持};
\node[box] (spice) at (-4.3,1.35) {SPICE / DE 历表};
\node[box] (r2s2) at (-4.3,0) {r2s2\\时空转换};
\node[box] (gmat2) at (4.3,1.35) {GMAT\\校核 / 运营分析};
\node[box] (cesium) at (4.3,0) {Cesium\\可视化（CZML）};
\node[box] (orekit) at (4.3,-1.35) {Orekit\\地球段 / 定轨};
\node[box] (pykep) at (-4.3,-1.35) {pykep + pygmo\\低推力全局优化};
\draw[arr] (spice.east) -- (e2.west |- spice);
\draw[arr] (r2s2.east) -- (e2.west |- r2s2);
\draw[arr] (e2.east |- gmat2) -- (gmat2.west);
\draw[arr] (e2.east) -- (cesium.west);
\draw[arr] (e2.east |- orekit) -- (orekit.west);
\draw[arr] (e2.west |- pykep) -- (pykep.east);
\draw[arr] (pykep.north) -- ++(0,0.4) -| (e2.south) node[pos=0.3, below, font=\scriptsize, text=e2gray] {优化解回代验证};
\end{tikzpicture}
\caption{e2m2e 的组合生态：上游时空/星历供给，下游可视化、校核、
定轨与全局优化各有现成开源件，e2m2e 专守地月设计引擎一层。}
\label{fig:eco}
\end{figure}

\begin{keybox}{小结：一个被反复验证的空位}
通用平台没有三体内核（GMAT）、有框架没有自动化（STK/Astrogator）、
有构件没有流水线（Orekit）、有动力学没有设计工具（pykep）、教学库
归档在 contrib 草稿（poliastro）——八个对象无一同时具备``八族生成 +
流形低能转移 + 三体保持 + 星历修正链''。这条被代码级调研逐一证实的
空位，就是 e2m2e 的生态位。
\end{keybox}
